% GENERATED FILE. Edit canonical lessons, then run npm run generate:books.
% canonical-source-hash: fnv1a64:a9517f74d29a8fcf
% canonical-lessons: BN-C428-ardro, BN-C428-bhepsa, BN-C428-jhapsa, BN-C428-setsete, BN-C428-dheke-jaoya

\chapter{Damp and Haze}
\label{ch:bn-damp-and-haze}
\hlchaptermodality{428}

\begin{chapteropening}
\textbf{By the end of this chapter:} \emph{I can say the air is humid and muggy, a wall is damp and the hills are hazy.}

Complete the last lesson of chapter 428: I can say the air is humid and muggy, a wall is damp and the hills are hazy.
\end{chapteropening}

\section[\texorpdfstring{\bn{আর্দ্র} (ārdrô)}{আর্দ্র (ārdrô)}]{\bn{আর্দ্র} (ārdrô) — humid; moist}

\label{lesson:BN-C428-ardro}



\begin{quote}
Before the new one: say the Bengali for temperature, then the Bengali for heat (warmth).
\end{quote}

\subsection*{You'll want to know: \bn{আর্দ্র}}
\textbf{\bn{আর্দ্র}} — \emph{ārdrô} — \textquotedblleft{}humid; moist\textquotedblright{}.

\begin{quote}
\bn{সমুদ্রের} \bn{কাছে} \bn{বাতাস} \bn{আর্দ্র।}
\end{quote}

\begin{quote}
\emph{sômudrer kāchhe bātās ārdrô.}
\end{quote}

\begin{quote}
Near the sea the air is humid.
\end{quote}

A careful word, heard in weather reports.

\begin{cousinweb}
\textbf{\bn{আর্দ্র}} (\emph{ārdrô}) is Sanskrit \emph{ārdra}, \textquotedblleft{}wet, moist, damp\textquotedblright{} (Monier-Williams), taken over whole.
\end{cousinweb}

\begin{grammarlens}[title={What you've built}]
One more word for getting around.
\end{grammarlens}

\subsection*{Your turn}
\begin{itemize}
\raggedright
  \item \emph{Say it:} \emph{ārdrô}
  \item \emph{Say it:} \emph{ārdrô}, once more
  \item \emph{Say it:} \emph{sômudrer kāchhe bātās ārdrô.} — Near the sea the air is humid.
  \item \emph{Recall:} say the Bengali for temperature, then the Bengali for heat (warmth), then say \emph{ārdrô} again
\end{itemize}

\begin{tcolorbox}[breakable,colback=teal!4,colframe=teal!35!black,title={Before you move on}]
What does \bn{আর্দ্র} mean? (\textquotedblleft{}Humid; moist\textquotedblright{}.) Say it once more.
\end{tcolorbox}
\section[\texorpdfstring{\bn{ভ্যাপসা} (bhæpsā)}{ভ্যাপসা (bhæpsā)}]{\bn{ভ্যাপসা} (bhæpsā) — muggy; close and sticky}

\label{lesson:BN-C428-bhepsa}



\begin{quote}
Before the new one: say the Bengali for heat (warmth), then the Bengali for humid.
\end{quote}

\subsection*{You'll want to know: \bn{ভ্যাপসা}}
\textbf{\bn{ভ্যাপসা}} — \emph{bhæpsā} — \textquotedblleft{}muggy; close and sticky\textquotedblright{}.

\begin{quote}
\bn{বৃষ্টির} \bn{আগে} \bn{ভ্যাপসা} \bn{গরম} \bn{ছিল।}
\end{quote}

\begin{quote}
\emph{brishṭir āge bhæpsā gôrôm chhilô.}
\end{quote}

\begin{quote}
Before the rain it was muggy and hot.
\end{quote}

Also of a smell: musty, stale.

\begin{cousinweb}
\textbf{\bn{ভ্যাপসা}} (\emph{bhæpsā}) is an expressive word for close, damp heat; its origin is not settled.
\end{cousinweb}

\begin{grammarlens}[title={What you've built}]
One more word for getting around.
\end{grammarlens}

\subsection*{Your turn}
\begin{itemize}
\raggedright
  \item \emph{Say it:} \emph{bhæpsā}
  \item \emph{Say it:} \emph{bhæpsā}, once more
  \item \emph{Say it:} \emph{brishṭir āge bhæpsā gôrôm chhilô.} — Before the rain it was muggy and hot.
  \item \emph{Recall:} say the Bengali for heat (warmth), then the Bengali for humid, then say \emph{bhæpsā} again
\end{itemize}

\begin{tcolorbox}[breakable,colback=teal!4,colframe=teal!35!black,title={Before you move on}]
What does \bn{ভ্যাপসা} mean? (\textquotedblleft{}Muggy; close and sticky\textquotedblright{}.) Say it once more.
\end{tcolorbox}
\section[\texorpdfstring{\bn{ঝাপসা} (jhāpsā)}{ঝাপসা (jhāpsā)}]{\bn{ঝাপসা} (jhāpsā) — hazy; blurred}

\label{lesson:BN-C428-jhapsa}



\begin{quote}
Before the new one: say the Bengali for humid, then the Bengali for muggy.
\end{quote}

\subsection*{You'll want to know: \bn{ঝাপসা}}
\textbf{\bn{ঝাপসা}} — \emph{jhāpsā} — \textquotedblleft{}hazy; blurred\textquotedblright{}.

\begin{quote}
\bn{কুয়াশায়} \bn{দূরের} \bn{পাহাড়} \bn{ঝাপসা} \bn{দেখাল।}
\end{quote}

\begin{quote}
\emph{kuyāshāy dūrer pāhāṛ jhāpsā dekhālô.}
\end{quote}

\begin{quote}
In the fog the distant hills looked hazy.
\end{quote}

Also of sight: \bn{চোখে} \bn{ঝাপসা} \bn{দেখা}, to see things blurred.

\begin{cousinweb}
\textbf{\bn{ঝাপসা}} (\emph{jhāpsā}) is an expressive word for what cannot be seen clearly; its origin is not settled.
\end{cousinweb}

\begin{grammarlens}[title={What you've built}]
One more word for getting around.
\end{grammarlens}

\subsection*{Your turn}
\begin{itemize}
\raggedright
  \item \emph{Say it:} \emph{jhāpsā}
  \item \emph{Say it:} \emph{jhāpsā}, once more
  \item \emph{Say it:} \emph{kuyāshāy dūrer pāhāṛ jhāpsā dekhālô.} — In the fog the distant hills looked hazy.
  \item \emph{Recall:} say the Bengali for humid, then the Bengali for muggy, then say \emph{jhāpsā} again
\end{itemize}

\begin{tcolorbox}[breakable,colback=teal!4,colframe=teal!35!black,title={Before you move on}]
What does \bn{ঝাপসা} mean? (\textquotedblleft{}Hazy; blurred\textquotedblright{}.) Say it once more.
\end{tcolorbox}
\section[\texorpdfstring{\bn{স্যাঁতসেঁতে} (sæ̃tsẽte)}{স্যাঁতসেঁতে (sæ̃tsẽte)}]{\bn{স্যাঁতসেঁতে} (sæ̃tsẽte) — damp; dank}

\label{lesson:BN-C428-setsete}



\begin{quote}
Before the new one: say the Bengali for muggy, then the Bengali for hazy.
\end{quote}

\subsection*{You'll want to know: \bn{স্যাঁতসেঁতে}}
\textbf{\bn{স্যাঁতসেঁতে}} — \emph{sæ̃tsẽte} — \textquotedblleft{}damp; dank\textquotedblright{}.

\begin{quote}
\bn{বর্ষায়} \bn{ঘরের} \bn{দেয়াল} \bn{স্যাঁতসেঁতে} \bn{হয়।}
\end{quote}

\begin{quote}
\emph{bôrshāy ghôrer deyāl sæ̃tsẽte hôy.}
\end{quote}

\begin{quote}
In the monsoon the walls of the house get damp.
\end{quote}

\begin{cousinweb}
\textbf{\bn{স্যাঁতসেঁতে}} (\emph{sæ̃tsẽte}) is inherited: CDIAL lists Bengali \emph{sẽt-sẽtiyā}, \textquotedblleft{}damp, moist\textquotedblright{}, under the reconstructed \emph{saṁtimita} (CDIAL 12891); the doubled shape is an expressive repetition.
\end{cousinweb}

\begin{grammarlens}[title={What you've built}]
One more word for getting around.
\end{grammarlens}

\subsection*{Your turn}
\begin{itemize}
\raggedright
  \item \emph{Say it:} \emph{sæ̃tsẽte}
  \item \emph{Say it:} \emph{sæ̃tsẽte}, once more
  \item \emph{Say it:} \emph{bôrshāy ghôrer deyāl sæ̃tsẽte hôy.} — In the monsoon the walls of the house get damp.
  \item \emph{Recall:} say the Bengali for muggy, then the Bengali for hazy, then say \emph{sæ̃tsẽte} again
\end{itemize}

\begin{tcolorbox}[breakable,colback=teal!4,colframe=teal!35!black,title={Before you move on}]
What does \bn{স্যাঁতসেঁতে} mean? (\textquotedblleft{}Damp; dank\textquotedblright{}.) Say it once more.
\end{tcolorbox}
\section[\texorpdfstring{\bn{ঢেকে} \bn{যাওয়া} (ḍheke jāoyā)}{ঢেকে যাওয়া (ḍheke jāoyā)}]{\bn{ঢেকে} \bn{যাওয়া} (ḍheke jāoyā) — to be covered over; to be shrouded}

\label{lesson:BN-C428-dheke-jaoya}



\begin{quote}
Before the new one: say the Bengali for hazy, then the Bengali for damp (and dank).
\end{quote}

\subsection*{You'll want to know: \bn{ঢেকে} \bn{যাওয়া}}
\textbf{\bn{ঢেকে} \bn{যাওয়া}} — \emph{ḍheke jāoyā} — \textquotedblleft{}to be covered over; to be shrouded\textquotedblright{}.

\begin{quote}
\bn{কুয়াশায়} \bn{পুরো} \bn{গ্রামটা} \bn{ঢেকে} \bn{গেল।}
\end{quote}

\begin{quote}
\emph{kuyāshāy puro grāmṭā ḍheke gelô.}
\end{quote}

\begin{quote}
The whole village was shrouded in fog.
\end{quote}

Past (he/she …): \bn{সেই} \bn{পাহাড়} \bn{মেঘে} \bn{ঢেকে} \bn{গেল}, \emph{sei pāhāṛ meghe ḍheke gelô}.

\bn{ঢেকে} is the -\bn{এ} form of \bn{ঢাকা}, to cover: what covers takes -\bn{এ}, \bn{কুয়াশায়}, \bn{মেঘে}.

\begin{cousinweb}
\textbf{\bn{ঢেকে} \bn{যাওয়া}} (\emph{ḍheke jāoyā}) joins \bn{ঢেকে}, the -\bn{এ} form of \bn{ঢাকা}, \textquotedblleft{}to cover\textquotedblright{}, which CDIAL lists as Bengali \emph{ḍhākā} under the reconstructed \emph{ḍhakk}, \textquotedblleft{}cover\textquotedblright{} (CDIAL 5574), to \bn{যাওয়া}, to go: to go under a cover.
\end{cousinweb}

\begin{grammarlens}[title={What you've built}]
One more word for getting around.
\end{grammarlens}

\subsection*{Your turn}
\begin{itemize}
\raggedright
  \item \emph{Say it:} \emph{ḍheke jāoyā}
  \item \emph{Say it:} \emph{ḍheke jāoyā}, once more
  \item \emph{Say it:} \emph{kuyāshāy puro grāmṭā ḍheke gelô.} — The whole village was shrouded in fog.
  \item \emph{Recall:} say the Bengali for hazy, then the Bengali for damp (and dank), then say \emph{ḍheke jāoyā} again
\end{itemize}

\begin{tcolorbox}[breakable,colback=teal!4,colframe=teal!35!black,title={Before you move on}]
What does \bn{ঢেকে} \bn{যাওয়া} mean? (\textquotedblleft{}To be covered over; to be shrouded\textquotedblright{}.) Say it once more.
\end{tcolorbox}
